\documentclass[10pt,a4paper]{amsart}
\usepackage[margin=19mm]{geometry}
\usepackage{amsmath,amssymb,mathtools}
\usepackage[scr=rsfs,cal=euler]{mathalfa}
\usepackage{xfrac}
\usepackage[unicode,hidelinks,bookmarks=false]{hyperref}
\newcommand{\ie}{i.e.\ }
\newcommand{\cf}{cf.\ }
\newcommand{\resp}{respectively\ }
\newcommand{\Subsection}{\subsection}
\providecommand{\nobreakdash}{\nobreak\hbox{-}}
\setlength{\emergencystretch}{2em}
\multlinegap0pt
\usepackage{xcolor}
\usepackage{enumerate}
\usepackage[shortlabels]{enumitem}
\usepackage{tikz}
\usepackage{amssymb}
\newtheorem{lemma}{Lemma}
\title{Convergence of a Critical Multitype Bellman--Harris Process with Finite-Mean Lifetimes}
\author{E.T. Kolkovska, J.A. L\'{o}pez-Mimbela, J.H. Ram\'{\i}rez-Gonz\'{a}lez}
\date{Source: arXiv:2410.17369v3 (CC BY 4.0)}
\begin{document}
\maketitle

\noindent\emph{Source and licence.} Convergence of a Critical Multitype Bellman--Harris Process with Finite-Mean Lifetimes. Version arXiv:2410.17369v3 (CC BY 4.0), distributed by arXiv under the Creative Commons Attribution 4.0 International licence, http://creativecommons.org/licenses/by/4.0/, as recorded in the arXiv licence metadata for this exact version and shown on the abstract page https://arxiv.org/abs/2410.17369v3. Full licence notice: \texttt{licenses/arxiv-2410.17369-CC-BY-4.0.txt}.\par\smallskip

\noindent\textit{Editorial scope.} Counted unit: the article's lemma lemma3.8 (source0.tex lines 3010--3037). The article's complete original proof of it is delivered with it (source0.tex lines 3039--3167, one \texttt{proof} environment of 129 lines). No mathematics is written, repaired or re-derived: the statement and the proof are the article's own lines, in the article's own notation, and no retained line is substituted or re-flowed.  Retained uncounted as the source-local context the proof needs: lines 654--658; lines 2875--2896; lines 2931--2950.  Nothing was omitted from any retained span, so the package carries no omission record.  No retained line contains a citation, so the wrapper carries no bibliography and no bibliography entry.  No retained line contains a figure environment, an image-inclusion command or drawing code, so no image is delivered and the package declares no asset dependency.  Editorial formatting only, in addition: an AMS \texttt{amsart} preamble that restates the article's own declarations and macros used by the retained lines, namely the environments \texttt{lemma} on separate counters, together with the standard packages those lines need.  The retained environments print in this document as Lemma 1 in order of appearance, whereas the article prints the same items with the numbering of its own full document.  The complete native source of the article as distributed in the arXiv e-print for this exact version is bundled unchanged as \texttt{sources/167271/source0.tex}.
\begin{equation}\label{tail2}
F_i(0)=0,\qquad
F_i(t)<1 \text{ for all } t\ge0,\qquad
1-F_i(t)\le A t^{-\eta^{(i)}} \text{ for all } t>0,
\end{equation}

\begin{lemma}\label{lemaaux1}
Let \(T^{(1)},\dots,T^{(K)}\) be independent nonnegative random variables, where
\(T^{(i)}\) has distribution \(F_i\). Assume that
\[
1-F_i(t)\le A t^{-\eta^{(i)}},
\qquad t>0,\quad i=1,\dots,K,
\]
for some \(A>0\) and \(\eta^{(i)}>1\). Set
\[
\eta:=\min_{1\le i\le K}\eta^{(i)}
\]
and
\[
T^*:=\max_{1\le i\le K}T^{(i)}.
\]
Then there exists a constant \(\bar A>0\) such that
\begin{equation}\label{eq:max-tail-K}
P(T^*>t)\le \bar A t^{-\eta},
\qquad t>0.
\end{equation}
In particular, \(E(T^*)<\infty\).
\end{lemma}

\begin{lemma}\label{lemaaux2}
Let \(Y_1,Y_2,\dots\) be i.i.d. nonnegative random variables with mean
\(m<\infty\). Assume that
\[
P(Y_1>t)\le A t^{-\eta},
\qquad t>0,
\]
for some \(A>0\) and \(\eta>1\). Then, for every \(c>0\), there exists
\(C>0\) such that, for all sufficiently large \(n\) and all \(x\ge cn\),
\begin{equation}\label{cotaeta1}
P\left(\sum_{k=1}^nY_k-nm\ge x\right)
\le C n x^{-\eta}.
\end{equation}
Consequently, for every \(\delta>0\), there exists \(C_\delta>0\) such that,
for all sufficiently large \(n\),
\begin{equation}\label{cotaeta2}
P\left(\frac1n\sum_{k=1}^nY_k-m\ge\delta\right)
\le C_\delta n^{1-\eta}.
\end{equation}
\end{lemma}

\begin{lemma}\label{lemma3.8}
Assume \eqref{tail2}. Let \(Z=(Z_t)_{t\ge0}\) be the Markov renewal process on
\(\mathbf K=\{1,\dots,K\}\), and let \(n_t\) be the number of renewals up to time \(t\).
Let \(T^{(1)},\dots,T^{(K)}\) be independent random variables such that \(T^{(r)}\) has
distribution \(F_r\), and define
\[
T^*:=\max_{1\le r\le K}T^{(r)}.
\]
Set
\[
\bar\mu:=E(T^*)<\infty.
\]
If \(c>0\) satisfies
\[
c<\frac1{\bar\mu},
\]
then, for every initial type \(i\in\mathbf K\), there exists a constant \(C>0\) such that,
for all sufficiently large \(t\),
\begin{equation}\label{cotan}
P_i(n_t\le ct)
\le
C t^{1-\eta},
\end{equation}
where
\[
\eta:=\min_{1\le r\le K}\eta^{(r)}.
\]
\end{lemma}

\begin{proof}
Let \((Y_n)_{n\ge0}\) be the embedded Markov chain of the Markov renewal
process. Thus \(Y_0\) is the initial type. Conditionally on
\((Y_n)_{n\ge0}\), let the holding times be independent random variables
\((\widetilde T_n)_{n\ge0}\), where \(\widetilde T_n\) has distribution
\(F_{Y_n}\). Define
\[
\sigma_0:=0,
\qquad
\sigma_n:=\sum_{r=0}^{n-1}\widetilde T_r,
\qquad n\ge1.
\]
Then
\[
n_t=\max\{n\ge0:\sigma_n\le t\}.
\]
Put
\[
n:=\lfloor ct\rfloor+1.
\]
Then
\[
\{n_t\le ct\}\subseteq \{\sigma_n>t\}.
\]

We now compare the holding times of the Markov renewal process with
independent copies of \(T^*\). Let \(F_*\) be the distribution function of
\[
T^*:=\max_{1\le r\le K}T^{(r)}.
\]
Since \(T^{(r)}\le_{\mathrm{st}}T^*\) for every \(r=1,\ldots,K\), we have
\[
F_r(x)\ge F_*(x),
\qquad x\ge0,\quad r=1,\ldots,K.
\]
Let \(U_0,U_1,\ldots\) be i.i.d. random variables uniformly distributed on
\((0,1)\), independent of the embedded chain \((Y_n)_{n\ge0}\). Using the
usual generalized inverse
\[
F^{-1}(u):=\inf\{x\ge0:F(x)\ge u\},
\qquad 0<u<1,
\]
define
\[
\widetilde T_n:=F_{Y_n}^{-1}(U_n),
\qquad
T_n^*:=F_*^{-1}(U_n),
\qquad n\ge0.
\]
Then, conditionally on \((Y_n)_{n\ge0}\), the variables
\((\widetilde T_n)_{n\ge0}\) are independent and \(\widetilde T_n\) has
distribution \(F_{Y_n}\). Hence this construction has the law of the holding
times of the Markov renewal process. Moreover, since
\(F_{Y_n}\ge F_*\), the monotonicity of generalized inverses gives
\[
\widetilde T_n\le T_n^*,
\qquad n\ge0,
\]
almost surely. The variables \(T_0^*,T_1^*,\ldots\) are i.i.d. with
distribution \(F_*\), that is, with the distribution of \(T^*\). Consequently,
for every \(n\ge1\),
\[
\sigma_n
=
\sum_{r=0}^{n-1}\widetilde T_r
\le
\sum_{r=0}^{n-1}T_r^*
=:\bar S_n
\qquad\text{a.s.}
\]
Therefore,
\[
P_i(n_t\le ct)
\le
P_i(\sigma_n>t)
\le
P(\bar S_n>t).
\]

Since \(n=\lfloor ct\rfloor+1\), we have \(n/t\to c\). Because
\(c<1/\bar\mu\),
\[
\frac{t}{n}\longrightarrow \frac1c>\bar\mu.
\]
Therefore, there exists \(\delta>0\) such that, for all sufficiently large
\(t\),
\[
t-n\bar\mu\ge \delta n.
\]
Thus,
\[
P(\bar S_n>t)
=
P(\bar S_n-n\bar\mu>t-n\bar\mu)
\le
P(\bar S_n-n\bar\mu\ge \delta n).
\]

By Lemma~\ref{lemaaux1}, the distribution of \(T^*\) satisfies
\[
P(T^*>x)\le \bar A x^{-\eta},
\qquad x>0.
\]
 Applying Lemma~\ref{lemaaux2} to the i.i.d. sequence
 \[
 Y_k:=T_{k-1}^*,
 \qquad k\ge1,
 \]
 with mean \(\bar\mu\), we obtain
 \[
 P(\bar S_n-n\bar\mu\ge \delta n)
 \le
 C n^{1-\eta}
 \]
 for all sufficiently large \(n\). Since \(n=\lfloor ct\rfloor+1\), it follows
that
\[
P(\bar S_n-n\bar\mu\ge \delta n)
\le
C' t^{1-\eta}
\]
for all sufficiently large \(t\). Combining the preceding estimates gives
\[
P_i(n_t\le ct)
\le
C' t^{1-\eta},
\]
which proves \eqref{cotan}.
\end{proof}

\end{document}
